%% F3 and F4. Data: the T1 runs, no extra run.
\begin{figure}[t]
  \centering
  \placeholderfig{\linewidth}{2.6cm}{\textbf{F3: AP against $\sigma$,} two panels, V2X-ViT on each dataset. Curves: perfect pose, uncorrected, + FreeAlign, + AlignFormer. Ego only as a dashed line.}
  \caption{\textbf{F3.} One host per dataset, fixed in advance.}
  \why{\textbf{Why this figure:} T1 in one glance for one host per dataset, fixed in advance so it is not picked after seeing results. \textbf{What we conclude.} Whether the corrected curve stays flat and close to perfect pose as noise grows, while the uncorrected one falls.}
\end{figure}
\begin{figure}[t]
  \centering
  \placeholderfig{\linewidth}{2.6cm}{\textbf{F4: BEV, one V2X-Real frame at $\sigma=1$.} V2X-ViT uncorrected, then + AlignFormer, ground truth overlaid. Frame picked by rule: median per-frame gain.}
  \caption{\textbf{F4.} What the correction changes for a host.}
  \why{\textbf{Why this figure:} it makes the numbers concrete, showing what a correction changes in a host's output. It proves nothing alone, so the frame is chosen by a fixed rule. \textbf{What we conclude.} What misaligned detections look like, and that correction removes them.}
\end{figure}
